\documentclass[aspectratio=169]{beamer}

\usepackage[T1]{fontenc}
\usepackage[utf8]{inputenc}
\usepackage{booktabs}
\usepackage{array}
\usepackage{tabularx}
\usepackage{graphicx}
\usepackage{xcolor}

\definecolor{evblue}{HTML}{1F4E79}
\definecolor{evgreen}{HTML}{2E7D32}
\definecolor{evorange}{HTML}{C55A11}
\definecolor{evgray}{HTML}{5F6B73}
\definecolor{evlight}{HTML}{F3F6F8}

\setbeamertemplate{navigation symbols}{}
\setbeamercolor{title}{fg=evblue}
\setbeamercolor{frametitle}{fg=evblue}
\setbeamercolor{structure}{fg=evblue}
\setbeamercolor{block title}{fg=white,bg=evblue}
\setbeamercolor{block body}{fg=black,bg=evlight}
\setbeamerfont{frametitle}{series=\bfseries}

\newcolumntype{Y}{>{\raggedright\arraybackslash}X}

\newcommand{\kw}{\mathrm{kW}}
\newcommand{\mw}{\mathrm{MW}}
\newcommand{\kwh}{\mathrm{kWh}}
\newcommand{\mwh}{\mathrm{MWh}}
\newcommand{\gwh}{\mathrm{GWh}}
\newcommand{\takeaway}[1]{\vspace{0.5em}\textcolor{evblue}{\textbf{Takeaway:}} #1}
\newcommand{\caveat}[1]{\textcolor{evorange}{\textbf{Caveat:}} #1}
\newcommand{\maybegraphic}[2][]{%
  \IfFileExists{#2}{%
    \includegraphics[#1]{#2}%
  }{%
    \setlength{\fboxsep}{1em}%
    \fbox{%
      \begin{minipage}[c][0.42\textheight][c]{0.82\linewidth}%
        \centering
        \textcolor{evgray}{\textbf{Figure not available}}\\[0.4em]
        \footnotesize\texttt{\detokenize{#2}}%
      \end{minipage}%
    }%
  }%
}

\title[Albany EV Charging Results]{Albany County Behavioral EV Charging Results}
\subtitle{Full dataset run with ZIP/ZCTA adoption scaling and H3-8 charging geography}
\author{NYSERDA EV Adoption and Charging Model}
\date{June 16, 2026}

\begin{document}

% -----------------------------------------------------------------------------
\begin{frame}
  \titlepage
\end{frame}

\begin{frame}{This week's status}
\small
\begin{tabularx}{\linewidth}{p{0.28\linewidth}Y}
\toprule
\textbf{Area} & \textbf{Current status} \\
\midrule
Model geography & Adoption now scales from home ZIP/ZCTA; charging load is placed at destination H3. \\
Charging behavior & Allocation now uses conditional expected value instead of simple expected-capacity scaling. \\
Managed charging & Applied only to charger types marked as managed-eligible in assumptions. \\
\bottomrule
\end{tabularx}
\end{frame}

% -----------------------------------------------------------------------------
\begin{frame}{Remaining trip-chain modeling decision: final stops}
\small
\textbf{Final activity distribution for chains that do not end at home}

\vspace{0.4em}
\begin{tabular}{lr}
\toprule
\textbf{Final activity type} & \textbf{Chains} \\
\midrule
Social / Recreational & 7,572 \\
Work & 3,300 \\
School / Daycare / Religious & 1,620 \\
Shopping / Errands & 909 \\
Other & 715 \\
Transport / drop-off / pick-up & 450 \\
Meals & 117 \\
Medical / Dental & 105 \\
\bottomrule
\end{tabular}

\begin{itemize}
  \item \textbf{As observed:} final non-home stops remain non-home overnight stops.
  \item \textbf{Completed chain:} add a synthetic return-home leg and overnight home dwell.
\end{itemize}
\end{frame}

% -----------------------------------------------------------------------------
\begin{frame}{Geography design: home adoption vs. physical load}
\small
\begin{block}{Core decision}
Keep home adoption geography separate from charging-load geography.
\end{block}

\begin{tabularx}{\linewidth}{p{0.26\linewidth}p{0.26\linewidth}Y}
\toprule
\textbf{Use case} & \textbf{Target geography} & \textbf{Reason} \\
\midrule
EV adoption forecast & Home ZIP/ZCTA & Adoption model is ZIP-level; join from home coordinate to ZIP/ZCTA. \\
Physical charging load & Destination H3 & Load appears where vehicles charge, not necessarily where they live. \\
\bottomrule
\end{tabularx}
\end{frame}


% -----------------------------------------------------------------------------
\begin{frame}{Charging allocation update}
\footnotesize
\setlength{\abovedisplayskip}{0.25em}
\setlength{\belowdisplayskip}{0.25em}
\begin{block}{Old risk}
Expected capacity logic could under- or over-allocate energy when a lower-priority stop had partial capacity and uncertain availability.
\end{block}

\begin{block}{Current method}
Carry probability states across ranked charging opportunities:
\[
\mathcal{S}_j = \{(p_k, r_k): \text{probability } p_k, \text{ remaining energy } r_k\}
\]
At each stop, split every state into available and unavailable branches, allocate:
\[
E_{j,k}=\min(r_k, C_j)
\]
and aggregate expected event energy:
\[
\mathbb{E}[E_j]=\sum_k p_k \cdot q_j \cdot E_{j,k}
\]
where $C_j$ is stop capacity and $q_j$ is the GridUp-style charger availability:
\[
q_j=\min\left(1,\ q^{\mathrm{peak}}_j\cdot
\frac{\mathrm{peak\ stopped\ vehicles}}{\mathrm{current\ stopped\ vehicles}}\right).
\]
\end{block}
\end{frame}

% -----------------------------------------------------------------------------
\begin{frame}{Managed vs. unmanaged charging}
\footnotesize
\begin{columns}[T,totalwidth=\linewidth]
\begin{column}{0.48\linewidth}
\begin{block}{Unmanaged}
Charging begins at arrival and uses charger power until allocated energy is delivered or the dwell window closes:
\[
T_s = \min\left(\tau_s, \frac{E_s}{P_s}\right)
\]
\end{block}
\end{column}
\begin{column}{0.48\linewidth}
\begin{block}{Managed}
Only managed-eligible charger types are spread across the dwell window:
\[
\bar{P}_s = \min\left(P_s, \frac{E_s}{\tau_s}\right)
\]
\end{block}
\end{column}
\end{columns}

\vspace{0.3em}
\begin{itemize}
  \item Managed charging applies to:
  \begin{itemize}
    \item home L2: 7.2 kW, managed eligible;
    \item work L2: 7.2 kW, managed eligible
    \item civic, long public, and quick public: not managed eligible
  \end{itemize}
  \item Availability is now GridUp-style: SFH home L2 has 100\% peak-hour probability, MFH home L2 and work L2 use 50\%, then probabilities rise in lower-demand hours.
  \item This keeps home/work management separate from public/DCFC charging behavior.
\end{itemize}
\end{frame}

% -----------------------------------------------------------------------------
\begin{frame}{Adoption scaling logic}
\small
\begin{block}{Forecast scaling equation}
\[
L_{g,t,c,y} = \sum_h L^{\text{fullEV}}_{h,g,t,c}
\cdot A_{h,y,s} \cdot G_{h,y,s}
\]
\end{block}

\begin{tabularx}{\linewidth}{p{0.16\linewidth}Y}
\toprule
\textbf{Symbol} & \textbf{Meaning} \\
\midrule
$h$ & Home ZIP/ZCTA used for adoption attribution. \\
$g$ & Charging H3 cell where load physically appears. \\
$c$ & Charger category / charging location type. \\
$A_{h,y,s}$ & Adoption fraction for home geography, year, and scenario. \\
$G_{h,y,s}$ & Vehicle growth factor for home geography, year, and scenario. \\
\bottomrule
\end{tabularx}
\end{frame}

% -----------------------------------------------------------------------------
\begin{frame}{Result Summary}
\small
\begin{tabular}{p{0.37\linewidth}p{0.52\linewidth}}
\toprule
\textbf{Question} & \textbf{Current result} \\
\midrule
Full-EV daily energy delivered & 3.305 \(\gwh\) per modeled day under both managed and unmanaged modes. \\
Full-EV peak load & 287.7 \(\mw\) unmanaged at 20:00; 197.8 \(\mw\) managed at 03:00. \\
2027 baseline adoption load & 155.0 \(\mwh\) per modeled day; 13.49 \(\mw\) unmanaged peak. \\
Managed charging effect & 31.3\% lower full-EV peak and 31.3\% lower 2027 baseline peak. \\
\bottomrule
\end{tabular}

\takeaway{The model now produces a home-ZCTA by charging-H3 by time load surface, then scales it with the ZIP-level adoption forecast.}
\end{frame}

% -----------------------------------------------------------------------------
\begin{frame}{Modeling Geography}
\small
\begin{columns}[T,totalwidth=\linewidth]
\begin{column}{0.49\linewidth}
\textbf{Adoption scaling geography}
\begin{itemize}
  \item Adoption forecast is ZIP/ZCTA-level.
  \item Each modeled vehicle day keeps its home ZCTA.
  \item Scaling is applied by \texttt{home\_zcta}, forecast year, and adoption scenario.
\end{itemize}
\end{column}
\begin{column}{0.49\linewidth}
\textbf{Charging-load geography}
\begin{itemize}
  \item Charging load is placed at the destination charging location.
  \item Primary planning geography is H3 at resolution 8.
\end{itemize}
\end{column}
\end{columns}
\end{frame}

% % -----------------------------------------------------------------------------
% \begin{frame}{Output Tables Produced}
% \scriptsize
% \renewcommand{\arraystretch}{1.15}
% \begin{tabular}{p{0.72\linewidth}r}
% \toprule
% \textbf{What the table contains} & \textbf{Rows} \\
% \midrule
% Expected charging sessions: timing, energy, charger type, event probability, and home/destination geography. & 1,445,636 \\
% Vehicle-day energy accounting by managed/unmanaged mode: daily energy need, delivered energy, and unmet energy. & 433,920 \\
% Pre-adoption load surface by home ZCTA, charging H3, charger type, and 15-minute time bin. & 2,437,544 \\
% Hourly version of the pre-adoption home-ZCTA by charging-H3 load surface. & 774,262 \\
% Full-EV charging-H3 load rollups used for spatial peak and profile analysis. & 644,185 / 176,766 \\
% Destination-point load rollups retained for QA and possible future feeder or transformer joins. & 8,884,374 / 2,479,365 \\
% Adoption-scaled charging-H3 planning loads across forecast years and adoption scenarios. & 19,325,550 / 5,302,980 \\
% \bottomrule
% \end{tabular}
% \end{frame}

% -----------------------------------------------------------------------------
\begin{frame}{Full-EV Energy Accounting}
\small
\begin{tabular}{p{0.43\linewidth}r}
\toprule
\textbf{Metric} & \textbf{Value per modeled day} \\
\midrule
Vehicle days & 216,960 \\
Total network miles & 8.964 million miles \\
Mean daily network miles & 41.3 miles per vehicle day \\
Daily travel energy need & 3.765 \(\gwh\) \\
Expected charging energy delivered & 3.305 \(\gwh\) \\
Unmet energy & 0.460 \(\gwh\) \\
Unmet share of daily travel energy & 12.2\% \\
\bottomrule
\end{tabular}

\vspace{0.8em}
\begin{block}{Interpretation}
The unmet-energy metric is not a grid load. It is a behavioral feasibility signal showing where the current stop, dwell, charger-access, and power assumptions cannot fully replenish modeled daily travel energy.
\end{block}
\end{frame}

% -----------------------------------------------------------------------------
\begin{frame}{Full-EV Hourly Load Shape}
\begin{center}
\maybegraphic[width=0.86\linewidth]{figures/full_ev_hourly_profile.pdf}
\end{center}

\small
\takeaway{Managed charging shifts the aggregate peak from 20:00 to 03:00 and reduces the full-EV peak from 287.7 \(\mw\) to 197.8 \(\mw\).}
\end{frame}

% -----------------------------------------------------------------------------
\begin{frame}{Charging Energy Is Concentrated at Home}
\begin{center}
\maybegraphic[width=0.74\linewidth]{figures/energy_share_2027_baseline.pdf}
\end{center}

\small
\takeaway{Home charging accounts for 87.5\% of 2027 baseline daily energy; public, work, and civic charging together account for the remaining 12.5\%.}
\end{frame}

% -----------------------------------------------------------------------------
\begin{frame}{2027 Baseline Adoption-Scaled Load}
\begin{center}
\maybegraphic[width=0.86\linewidth]{figures/baseline_2027_hourly_profile.pdf}
\end{center}

\small
\takeaway{For 2027 baseline adoption, managed charging lowers the county-wide hourly peak from 13.49 \(\mw\) to 9.27 \(\mw\), with the same total daily energy of 155.0 \(\mwh\).}
\end{frame}

% -----------------------------------------------------------------------------
\begin{frame}{Adoption-Scaled Energy by Forecast Year}
\begin{center}
\maybegraphic[width=0.86\linewidth]{figures/scaled_energy_by_year.pdf}
\end{center}
\end{frame}

% -----------------------------------------------------------------------------
\begin{frame}{Adoption-Scaled Peak by Forecast Year}
\begin{center}
\maybegraphic[width=0.86\linewidth]{figures/scaled_peak_by_year.pdf}
\end{center}
\end{frame}

% -----------------------------------------------------------------------------
\begin{frame}{2027 Baseline H3 Hotspots}
\small
\begin{columns}[T,totalwidth=\linewidth]
\begin{column}{0.49\linewidth}
\textbf{Unmanaged top H3-hours}

\vspace{0.3em}
\scriptsize
\begin{tabular}{lrr}
\toprule
\textbf{H3 cell} & \textbf{Hour} & \textbf{\(\kw\)} \\
\midrule
882b8902b1fffff & 20 & 162.7 \\
882b8902b1fffff & 19 & 157.4 \\
882b890297fffff & 20 & 155.2 \\
882b890025fffff & 19 & 154.2 \\
882b8902b1fffff & 21 & 153.3 \\
\bottomrule
\end{tabular}
\end{column}
\begin{column}{0.49\linewidth}
\textbf{Managed top H3-hours}

\vspace{0.3em}
\scriptsize
\begin{tabular}{lrr}
\toprule
\textbf{H3 cell} & \textbf{Hour} & \textbf{\(\kw\)} \\
\midrule
882b8902b1fffff & 4 & 108.7 \\
882b8902b1fffff & 3 & 108.7 \\
882b8902b1fffff & 2 & 107.0 \\
882b890297fffff & 3 & 106.5 \\
882b8902b1fffff & 1 & 106.4 \\
\bottomrule
\end{tabular}
\end{column}
\end{columns}
\end{frame}

% -----------------------------------------------------------------------------
\begin{frame}{Current Caveats}
\small
\begin{itemize}
  \item Charger-access and charger-power assumptions are configurable but not yet calibrated against local charger inventory and housing data.
  \item The adoption forecast in this run covers 2018--2027, so longer-term grid planning requires a longer adoption horizon.
\end{itemize}
\end{frame}

% -----------------------------------------------------------------------------
\begin{frame}{Next Steps}
\small
\begin{enumerate}
  \item Map charging H3 cells to feeders, substations, or utility service polygons when those layers become available.
  \item Calibrate charger availability by home type, workplace access, and public EVSE inventory.
  \item Extend the ZIP/ZCTA adoption forecast beyond 2027 before using the outputs for long-horizon planning.
  \item Run sensitivity cases for charger access, managed eligibility, DCFC share, and energy efficiency.
\end{enumerate}
\end{frame}

\end{document}
